\chapter{格式、错误码与接口索引}
本附录把各章反复使用的稳定约定集中列出，方便实现时查阅；若和练习里的公开签名不一致，以课程 manifest 和 Java 类型为准。本课程定义的布局与错误码只属于 simpleKafka 教学接口，除非明确引用官方资料，不应把它们称为 Kafka wire、存储格式或配置默认值。

\section{磁盘与网络格式}
\subsection{记录与索引}
记录为大端编码：
\begin{center}
\begin{tabular}{p{.22\textwidth}p{.17\textwidth}p{.53\textwidth}}
\toprule
字段 & 宽度 & 语义 \\
\midrule
length & int32 & 不含自身的记录长度，含 CRC；28 至 1,048,576。\\
crc32c & int32 & 校验 offset 到 value 末尾，不含 length/CRC 字段自身。\\
offset & int64 & 非负逻辑位点。\\
timestamp & int64 & 非负时间戳。\\
keyLength & int32 & -1 为 null，0 为存在的空 key，其余为 bytes 数。\\
valueLength & int32 & 非负 bytes 数，value 不为 null。\\
key/value & 变长 & 按前述长度紧跟在字段后。\\
\bottomrule
\end{tabular}
\end{center}
每条记录的固定长度部分（不含 length 前缀）包含 CRC、offset、timestamp、两个 length，共 28 bytes；长度前缀使最小总占用为 32 bytes。length 字段值下限仍是 28，因为其统计从 CRC 开始。每条 segment 的稀疏索引条目是 \pathcode{offset:int64,position:int64}，共 16 bytes；索引每 intervalRecords 条记录写一个点，首条总有点。索引可以丢弃并由日志重建。
普通记录保持原布局。stamped record 在 timestamp 后、keyLength 前插入以下 64-byte producer identity：\pathcode{marker:int32=-2 | producerId:int64 | producerEpoch:int32 | firstSequence:int64 | batchSize:int32 | batchIndex:int32 | batchHash:32 bytes}。完整布局为 \pathcode{length | crc32c | offset | timestamp | marker | producerId | producerEpoch | firstSequence | batchSize | batchIndex | batchHash | keyLength | valueLength | key | value}。最小总长为 96 bytes，1 MiB 上限包含 stamp。CRC32C 覆盖 offset 到最后一个 value byte，包含 stamp。普通 record 字节不变。
已提供的 \method{ProducerStamp.ENCODED_OVERHEAD} 为 64 bytes，磁盘与 \method{MessageCodec} wire record 的 1 MiB 上限均包含这部分开销。

目录约定为 \pathcode{<root>/<topic>/<partition>/<baseOffset:020d>.log}，配对索引后缀为 \pathcode{.index}。保留范围是 [logStartOffset,LEO)，fetch 左闭右开；LEO 是下一条可写 offset，HW 是普通 consumer 能观察的最高提交边界（仅 offset < HW 可见）。group committed offset 也是 next offset，但不等于 HW。

\subsection{TCP frame 与请求体}
教学 frame 为 \pathcode{length:int32 | version:int16(1) | api:int16 | correlationId:int32 | error:int16 | payload}，length 不含自身，范围 [10,8,388,608]。读到新 frame 首字节前 EOF 返回 null；frame 内 EOF 报 \method{INVALID_REQUEST}。请求 error 字段必须为 0。未知 version/API、过长 frame 和尾随 body 字节均拒绝。每条磁盘记录上限 1 MiB，不因 frame 上限而改变。

API id 固定如下：
\begin{center}
\begin{tabular}{rl|rl}
\toprule
ID & API & ID & API \\
\midrule
1 & METADATA & 6 & \method{JOIN_GROUP} \\
2 & PRODUCE & 7 & HEARTBEAT \\
3 & FETCH & 8 & \method{LEAVE_GROUP} \\
4 & \method{COMMIT_OFFSET} & 9 & \method{GROUP_ASSIGNMENT} \\
5 & \method{FETCH_OFFSET} & 10 & \method{REPLICA_FETCH} \\
11 & \method{IDEMPOTENT_PRODUCE} & & \\
\bottomrule
\end{tabular}
\end{center}
MessageCodec 使用大端基础字段；字符串为 int32 UTF-8 byte 长度后接 bytes，集合为 int32 数量后逐项，Optional 值用 boolean 再接值，byte[] 用 int32 长度（仅 key 可 -1）。请求/成功响应字段按 provided Messages nested record 声明顺序编码；map 结果按 key 排序。Reply 头部 error=NONE 时 body 是相应成功 record；错误时 body 为 ErrorBody(message)，错误码只在 frame header 出现。字段 schema 由 \pathcode{provided/.../protocol/Messages.java} 固定，学生不需重新定义十一种消息格式。FETCH 或 \pathcode{REPLICA_FETCH} wire 上的 stamped RecordData 为 \pathcode{marker=-2 | producerId | producerEpoch | firstSequence | batchSize | batchIndex | batchHash[32] | keyLength | key | valueLength | value | timestamp}；普通 record 的 wire 编码不变。
\subsection{消息 body 字段顺序}
下面的请求与响应由 MessageCodec 按表内参数声明顺序编码。集合长度和字符串/byte 长度均为大端 int32；列表项无额外 tag，按各 value type 声明顺序内联。GroupToken nullable，用 boolean present 后跟 group、member、generation；Acks 为 int16，LEADER=1、ALL=-1。
\setlength{\tabcolsep}{3pt}
\small
\begin{longtable}{>{\raggedright\arraybackslash}p{.22\textwidth}>{\raggedright\arraybackslash}p{.36\textwidth}>{\raggedright\arraybackslash}p{.36\textwidth}}
\toprule
API & 请求记录与字段 & 成功响应记录与字段 \\
\midrule
\endhead
\method{METADATA} & MetadataRequest: topic & MetadataBody: partitions \\
\method{PRODUCE} & ProduceRequest: tp, epoch, acks, timeoutMillis, records & ProduceBody: result \\
\method{FETCH} & FetchRequest: tp, epoch, offset, maxRecords, maxBytes, token & FetchBody: records, logStartOffset, logEndOffset, highWatermark, epoch \\
\method{IDEMPOTENT_PRODUCE} & IdempotentProduceRequest: tp, epoch, acks, timeoutMillis, producerId, producerEpoch, firstSequence, records & ProduceBody: result \\
\method{COMMIT_OFFSET} & CommitOffsetRequest: key, nextOffset, token & EmptyBody \\
\method{FETCH_OFFSET} & FetchOffsetRequest: key & OffsetBody: nextOffset (OptionalLong) \\
\method{JOIN_GROUP} & JoinGroupRequest: group, member, topic & GroupBody: assignment \\
\method{HEARTBEAT} & HeartbeatRequest: token & EmptyBody \\
\method{LEAVE_GROUP} & LeaveGroupRequest: group, member & GroupBody: assignment \\
\method{GROUP_ASSIGNMENT} & GroupAssignmentRequest: group, member & GroupBody: assignment \\
\method{REPLICA_FETCH} & ReplicaFetchRequest: tp, brokerId, epoch, fetchOffset, maxRecords, maxBytes, recoveryRead & ReplicaFetchBody: records, epoch, highWatermark, leaderLogEndOffset \\
\bottomrule
\end{longtable}
\normalsize
\begin{center}
\small
\begin{tabular}{>{\raggedright\arraybackslash}p{.25\textwidth}>{\raggedright\arraybackslash}p{.67\textwidth}}
\toprule
Record & Fields in declaration order \\
\midrule
\method{TopicPartition} & \method{topic}, \method{partition} \\
\method{RecordData} & \method{key}, \method{value}, \method{timestamp}, nullable \method{producerStamp} \\
\method{LogRecord} & \method{offset}, \method{data} \\
\method{AppendResult} & \method{firstOffset}, \method{nextOffset} \\
\method{OffsetKey} & \method{group}, \method{tp} \\
\method{Endpoint} & \method{host}, \method{port} \\
\method{GroupAssignment} & \method{generation}, \method{assignments} \\
\method{PartitionMetadata} & \method{tp}, \method{leaderId}, \method{epoch}, \method{replicas}, \method{isr}, \method{leader} \\
\method{ProducerStamp} & \method{producerId}, \method{producerEpoch}, \method{firstSequence}, \method{batchSize}, \method{batchIndex}, \method{batchHash} \\
\method{BusinessEvent} & \method{eventId}, \method{account}, \method{delta} \\
\bottomrule
\end{tabular}
\end{center}
\normalsize
Each byte array is length-prefixed; only key may be null (length -1), while value is always present but may be empty.

OffsetStore entries use \method{MessageCodec.encodeOffsetEntry}, with null key and timestamp 0. The record's physical log offset is independent of committed nextOffset.
BusinessEvent value 使用 \method{MessageCodec.encodeBusinessEvent} 编码：\pathcode{version:uint8=1 | eventId:identifier | account:identifier | delta:int64}。两个标识均为非空、严格 UTF-8 且最长 255 bytes；delta 可正、可负或为零。
\subsection{标识符与值类型}
topic、group、member 是非空 UTF-8 标识且最长 255 bytes；topic 只能包含 ASCII 字母数字、点、下划线、连字符，并拒绝 \pathcode{.} 与 \pathcode{..}。TopicPartition 的 partition 非负。RecordData 的 key 可以 null、value 不可 null（允许空数组），timestamp 非负；LogRecord offset 非负。数组所有权移交值对象后由接收方持有，不需要反复复制。所有 offset 表示下一条时都使用 half-open 边界，具体区别见正文各章。

\section{幂等追加与持久化副作用}
API id 11 是独立的 \method{IDEMPOTENT_PRODUCE} 请求；普通 \method{PRODUCE} 仍为非幂等模式，并拒绝带 stamp 的 client 输入。请求携带 \pathcode{tp, epoch, acks, timeoutMillis, producerId, producerEpoch, firstSequence, records}；broker 为每条记录加入请求的 batch identity，并对记录数量与 timestamp/key/value 内容流式计算 SHA-256 指纹。producer id 与 epoch 由调用方管理。新 producer epoch 从 sequence 0 开始，且只有记录持久化后才生效。完全相同的已完成请求返回原 offset 区间，不追加；相同 identity 携带不同内容会冲突；sequence gap 与未知的旧 sequence 会被拒绝。匹配的部分 batch 可只追加缺失后缀来完成。派生状态从保留日志历史恢复；当 \method{logStartOffset} 非零时，幂等请求报 \method{PRODUCER_STATE_EXPIRED}。本课程没有 producer-state snapshot 或 client session WAL。

\method{IdempotentProducer} 每次调用只发送一次。调用失败会保留 pending batch；只有调用方显式调用 \method{retry} 才刷新 metadata，并以同一 identity、sequence 和 payload 重发一次。重复请求仍通过普通 ACK precheck 并等待确认，包括 ALL 的 ISR/HW/deadline 条件。这不改变从不自动重试的 \method{SimpleProducer}。

\method{DurableEffectStore} 将一条 \method{BusinessEvent} 作为一条本地日志记录 force 写入。首次 event id 会追加并更新账户余额；id 与内容都相同的重放返回 false，不再追加；同 id 不同内容或 checked balance 溢出会被拒绝。恢复从保留日志重建事件身份和余额，并要求日志起点为零。该本地账本不与 \method{ProcessingLoop} 的 offset commit 原子协调：commit 失败后 callback 可再次运行，而稳定 event id 可防止本地记账重复。group fencing 不能撤销已经发生的副作用。
\section{固定错误码}
\begin{longtable}{r p{.27\textwidth} p{.55\textwidth}}
\toprule
Wire id & 名称 & 用途 \\
\midrule
\endhead
0 & NONE & 成功。\\
1 & \method{INVALID_REQUEST} & 参数、长度、版本或请求格式错误。\\
2 & \method{UNKNOWN_TOPIC_OR_PARTITION} & topic/partition 不存在。\\
3 & \method{OFFSET_OUT_OF_RANGE} & offset 超过当前日志起点/终点约束。\\
4 & \method{CORRUPT_RECORD} & 完整记录 CRC、格式或已提交历史冲突。\\
5 & \method{NOT_LEADER} & 请求者不是当前 leader。\\
6 & \method{FENCED_EPOCH} & 请求使用旧 leader epoch。\\
7 & \method{NOT_ENOUGH_REPLICAS} & 当前 ISR 少于 minISR。\\
8 & \method{REQUEST_TIMEOUT} & 请求 deadline 到期，结果可能未知。\\
9 & \method{ILLEGAL_GENERATION} & 消费组 generation 已过期。\\
10 & \method{UNKNOWN_MEMBER} & 消费组 member 未登记。\\
11 & \method{NOT_ASSIGNED} & 操作涉及未分配 partition。\\
12 & \method{NO_ELIGIBLE_LEADER} & 无满足干净选举条件的候选。\\
13 & \method{STORAGE_ERROR} & I/O 失败；保留底层 cause。\\
14 & \method{FENCED_PRODUCER_EPOCH} & 请求使用旧 producer epoch。\\
15 & \method{OUT_OF_ORDER_SEQUENCE} & batch sequence 或边界无效。\\
16 & \method{DUPLICATE_SEQUENCE_CONFLICT} & 重复 batch identity 的内容不同。\\
17 & \method{PRODUCER_STATE_EXPIRED} & 所需 producer 历史已不在保留日志前缀中。\\
18 & \method{INCOMPLETE_BATCH} & 无关追加会跨越未完成的 producer batch。\\
\bottomrule
\end{longtable}
未知 ID 不映射为成功；I/O 错误不吞 cause；网络错误 body 不回传服务器堆栈。

\section{三十一步与编辑入口}
所有路径均相对于仓库根目录，学生修改 \pathcode{exercises/src/main/java/io/simplekafka/} 下指定方法。每步累积命令为 \coursecmd{./gradlew :test -Pstep=N}；单步定位命令为 \coursecmd{./gradlew :stepTest -Pstep=N}，它不是晋级证明。
\begin{longtable}{r p{.20\textwidth} p{.68\textwidth}}
\toprule
步 & 课程标题 & 编辑方法 \\
\midrule
\endhead
1 & 记录编解码 & storage/RecordCodec: encode, decode \\
2 & 日志追加 & storage/SegmentLog: append \\
3 & 按位点读取 & storage/SegmentLog: readFrom \\
4 & 崩溃恢复 & storage/SegmentLog: recover \\
5 & 稀疏索引 & storage/SparseIndex: add, floor, rebuild \\
6 & 分段日志 & storage/PartitionLog: append, read \\
7 & 日志保留 & storage/PartitionLog: deleteBefore \\
8 & 尾部截断 & storage/PartitionLog: truncateTo \\
9 & 分区路由 & client/Partitioner: choose \\
10 & 主题元数据 & broker/PartitionCatalog: createTopic, metadata \\
11 & TCP帧协议 & protocol/FrameCodec: write, read \\
12 & 单机broker & broker/BrokerHandler: handle \\
13 & 批量生产者 & client/SimpleProducer: send, flush \\
14 & 拉取消费者 & client/SimpleConsumer: assign, seek, poll \\
15 & 消费位点持久化 & group/OffsetStore: commit, fetch;\newline client/SimpleConsumer: resume, commitSync;\newline broker/BrokerHandler: \method{COMMIT_OFFSET}, \method{FETCH_OFFSET} \\
16 & 至少一次处理 & client/ProcessingLoop: runOnce \\
17 & 分区分配 & group/RoundRobinAssignor: assign \\
18 & 成员与重平衡 & group/GroupCoordinator: join, leave, assignment;\newline broker/BrokerHandler: \method{JOIN_GROUP}, \method{LEAVE_GROUP};\newline \method{GROUP_ASSIGNMENT} \\
19 & 心跳与过期 & group/GroupCoordinator: heartbeat, expire;\newline broker/BrokerHandler: \method{HEARTBEAT} \\
20 & 消费组围栏 & group/GroupCoordinator: validate;\newline client/GroupConsumer: subscribe, poll, commitSync;\newline broker/BrokerHandler: \method{FETCH}/\method{COMMIT_OFFSET} token fencing \\
21 & 拉取复制 & broker/BrokerHandler: handle(\method{REPLICA_FETCH});\newline replication/FollowerReplicator: pollOnce \\
22 & ISR与高水位 & replication/ReplicationTracker:\newline report, expireLagging, advanceHighWatermark \\
23 & 生产确认 & replication/AckPolicy: validateBeforeAppend, await \\
24 & 并发副本broker & replication/ReplicatedPartition: produce, fetch \\
25 & 干净选举与epoch & replication/LeaderElection: elect, checkLeader \\
26 & 分叉日志修复 & replication/ReplicaReconciler: reconcile \\
27 & 元数据路由刷新 & client/MetadataRouter: refresh, call \\
28 & 副本重新入ISR & replication/ReplicaAdmission: tryAdd \\
29 & 三层一致性故障实验 & lab/ConsistencyExperiments: producerOutcomeUnknown, leaderAckLoss, allAckTimeout, consumerReplay, staleGroupSideEffect \\
30 & 幂等生产者 & client/IdempotentProducer: send, retry;\newline replication/IdempotentAppender: append, recover;\newline replication/ReplicatedPartition: produceIdempotent;\newline broker/BrokerHandler: handle(\method{IDEMPOTENT_PRODUCE});\newline client/MetadataRouter: call;\newline storage/RecordCodec: encode, decode;\newline storage/SegmentLog: append, recover;\newline storage/PartitionLog: append, deleteBefore, truncateTo \\
31 & 消费副作用幂等 & client/DurableEffectStore: apply, recover \\
\bottomrule
\end{longtable}
前七章各含四步；第 8 章含第 29–31 步。章节末累计通过点依次为 4、8、12、16、20、24、28、31。若前一步未通过，后一步测试通常也会失败，因为课程采用同一持续演进的系统而非互不相干的演示程序。

\section{课程简化与 Kafka 的差异}
\begin{enumerate}
\item \textbf{路由：}key 使用 CRC32C unsigned modulo；Kafka 默认 key partition 算法不同。null key 采用实例级 round-robin，而非 Kafka 的 sticky batching。
\item \textbf{磁盘：}记录逐条编码并每 batch 显式 \method{force(false)}，不是 Kafka record batch 与默认 page-cache 路径；强制写入不保证媒介级断电安全。
\item \textbf{索引：}每固定记录数写一个稀疏点，不按 bytes 间隔。
\item \textbf{消费组：}仅单 topic、即时重分配的简化模型，不是完整 classic barrier、cooperative 重平衡或 KIP-848。token 围栏仅保护课程 broker 内 offset 提交，不原子化外部副作用。
\item \textbf{offset：}本地 OffsetStore 不是高可用复制压缩的 \pathcode{__consumer_offsets}。
\item \textbf{复制与控制面：}各 broker 用真实 TCP 拉取并写独立目录；authority/coordinator 却在同 JVM 且单一可信，不是 KRaft，不保证 controller/authority 故障后的集群恢复或任意分区下共识。
\item \textbf{可见性假设：}本课程保留已提交前缀，要求干净候选 LEO 不低于 authority HW，使受控切换的已知提交边界不回退；这比对真实 Kafka 跨 epoch HW 行为更强，不可推广成 Kafka 永远单调的结论。
\item \textbf{重试、幂等与事务：}普通 \method{SimpleProducer} 从不自动重试，timeout 可能已追加。独立的 \method{IdempotentProducer} 支持由调用方管理 producer identity 并显式重试，batch stamp 持久化在日志中；但它不实现 Kafka PID 分配、transaction coordinator、client session WAL 或 producer-state snapshot，retention 可能使所需前缀过期。本地 \method{DurableEffectStore} 按稳定业务 event id 去重，但与 consumer offset commit 不原子协调。没有 Kafka 事务或 \method{read_committed} 保证、compaction、fire-and-forget 或 ELR。at-least-once 容许重复；本教学子集不声称端到端 exactly-once。
\item \textbf{恢复：}第 26 步全量比较前缀为 O(n) 教学简化；真实系统利用 leader epoch 历史缩小分叉搜索范围。
\end{enumerate}

\section{官方资料与延伸阅读}
以下资料用于解释 Kafka 概念与真实系统差异，并非课程自定义 wire/storage 布局的规范。版本标签反映本书所引用页面；不同版本配置和协议细节可能改变，阅读时以所列页面自身版本为准。
\begin{longtable}{p{.24\textwidth}p{.22\textwidth}p{.47\textwidth}}
\toprule
主题 & 官方资料 & 关联章节 \\
\midrule
\endhead
Kafka 设计与持久化模型 & \href{https://kafka.apache.org/42/design/design/}{Kafka 4.2 Design} & 日志、复制、消费者位置的官方概览；本课程简化项见上节。\\
消费位置与提交 & \href{https://kafka.apache.org/42/javadoc/org/apache/kafka/clients/consumer/KafkaConsumer.html}{KafkaConsumer 4.2 API} & 第 14–16 章：position、seek、commit 与线程约束。\\
最小 ISR 配置 & \href{https://kafka.apache.org/43/configuration/topic-configs/}{Kafka 4.3 Topic Configs} & 第 22–24 章：min.insync.replicas 与写入确认。\\
消费组协议差异 & \href{https://kafka.apache.org/43/operations/consumer-rebalance-protocol/}{Kafka 4.3 Rebalance Protocol} & 第 17–20 章：classic 与新协议背景；课程不实现完整协议。\\
leader epoch 与截断 & \href{https://cwiki.apache.org/confluence/display/KAFKA/KIP-101+-+Alter+Replication+Protocol+to+use+Leader+Epoch+rather+than+High+Watermark+for+Truncation}{KIP-101} & 第 26 章：为何仅靠 HW 不足以定位分叉。\\
JUnit 测试运行 & \href{https://docs.junit.org/5.11.4/user-guide/}{JUnit 5.11.4 User Guide} & 阅读实际 Step 行为失败与测试报告。\\
Tectonic 编译选项 & \href{https://tectonic-typesetting.github.io/book/latest/v2cli/compile.html}{Tectonic compile} & 本书采用 XeTeX/CJK 与本地 cache 的构建入口。\\
\bottomrule
\end{longtable}

\section{排版与构建说明}
默认主文件为英文版 \pathcode{simpleKafka.tex}，使用 \pathcode{book} 文档类；中文版 \pathcode{simpleKafka-zh.tex} 使用 \pathcode{ctexbook[UTF8,fontset=none,openany]}。共享 preamble 通过 fontspec 与 xeCJK 显式使用 TeX bundle 内的 Latin Modern 与 Fandol 字体文件，不探测系统字体；Noto、fontconfig、xelatex 与 latexmk 都不是构建依赖。代码片段使用 listings，不依赖 minted、外部绘图器、shell-escape 或联网网页嵌入。TikZ 图由 TeX 本地绘制。

\coursecmd{./gradlew :setup} 解析经过 checksum 校验的 JUnit；\coursecmd{./gradlew :setupBook} 为 PDF 准备 Tectonic。setup 不准备 PDF 引擎；\pathcode{:setupBook} 不编译 Java 或解析 JUnit，但 Wrapper 启动它仍需 Java 17+。\coursecmd{./gradlew :book} 默认生成英文版；\coursecmd{COURSE_LANG=zh ./gradlew :book} 生成中文版。英文 PDF 位于 \pathcode{build/book/en/simpleKafka.pdf}，中文 PDF 位于 \pathcode{build/book/zh/simpleKafka-zh.pdf}。\pathcode{COURSE_LANG} 只接受 \pathcode{en} 或 \pathcode{zh}。首次构建在线获取 TeX bundle 与字体，cache 位于仓库 \pathcode{.tools/tectonic-cache}；两种语言都在线成功后，可用 \pathcode{BOOK_OFFLINE=1} 启用 Tectonic only-cached 模式，Gradle \pathcode{--offline} 则独立控制 Gradle 依赖解析。PDF 本身不编译 Java，也不使用 JUnit。显式 \pathcode{TECTONIC} 可提供兼容引擎；\pathcode{TEX_BUNDLE} 接受 URL、本地目录、\pathcode{.zip} 或 \pathcode{.ttb} bundle 文件，本地普通 \pathcode{.tar} 不受 Tectonic 此接口支持。首次 clone 离线构建不受保证；bundle URL 也不是完整资源集的内容摘要，因此不承诺 PDF 字节级可重复。